\documentclass[10pt]{beamer}

% Theme selection
\usetheme{Madrid}
\usecolortheme{seahorse} 

% Required packages
\usepackage{amsmath, amssymb, amsfonts}
\usepackage{graphicx}
\usepackage{booktabs}
\usepackage{xcolor}
\usepackage{tikz}

% Xepersian must be the LAST package loaded
\usepackage{xepersian}

% Set the Persian font (Make sure this font is installed on your system or Overleaf)
\settextfont{Yas} % می‌توانید این را به B Nazanin یا Vazirmatn تغییر دهید

% Presentation Information
\title[کران‌های محرمانگی تفاضلی موضعی]{کران‌های آماری و نظریه اطلاعاتی محرمانگی تفاضلی موضعی}
\subtitle{(محرمانگی در یادگیری ماشین)}
\author[فیروزه ابریشمی]{فیروزه ابریشمی}
\institute[دانشگاه صنعتی شریف]{
	دانشکده علوم ریاضی\\
	دانشگاه صنعتی شریف\\
	\medskip
	\textit{استاد راهنما: آقای دکتر جواد ابراهیمی بروجنی}
}
\date{\today}

\setbeamersize{text margin left=1cm, text margin right=1.cm}

\begin{document}
	
	% اسلاید عنوان
	\begin{frame}
		\titlepage
	\end{frame}
	
	\begin{frame}{عصر کلان‌داده و ضرورت محرمانگی}
		\begin{itemize}
			\item \textbf{دنیای داده‌محور:} تصمیم‌گیری‌های مدرن و الگوریتم‌های یادگیری ماشین به شدت به مجموعه‌داده‌های عظیم (کلان‌داده) وابسته‌اند.
			
			\pause
			\vspace{0.25cm}
			\item \textbf{آسیب‌پذیری اصلی:} داده‌های خامی که برای آموزش استفاده می‌شوند، اغلب حاوی اطلاعات شخصی بسیار حساس (مانند سوابق پزشکی) هستند.
			
			\pause
			\vspace{0.25cm}
			\item \textbf{ناکارآمدی گمنام‌سازی سنتی:} صرفاً حذف شناسه‌های مستقیم (مانند نام‌ها و کدهای ملی) اساساً ناکافی است. 
			\begin{itemize}
				\item دانش جانبی امکان \textit{حملات شناسایی مجدد} فاجعه‌باری را فراهم می‌کند
				\\ (مانند داده‌های مسابقه نتفلیکس).
			\end{itemize}
			
			\pause
			\vspace{0.4cm}
			\item \textbf{هدف نهایی:} چگونه می‌توانیم بدون نقض محرمانگی افراد شرکت‌کننده، سودمندی آماری معناداری از داده‌ها استخراج کنیم؟
		\end{itemize}
	\end{frame}
	
	
	\begin{frame}{محرمانگی تفاضلی: شهود}
		\begin{itemize}
			\item \textbf{تضمین معنایی:} خروجی یک الگوریتم باید فارغ از حضور یا عدم حضور هر فرد خاصی در مجموعه‌داده، تقریباً یکسان بماند.
			\vspace{0.3cm}
			% \item \textbf{شهود هندسی (انقباض):} مکانیزم مانند فیلتری عمل می‌کند که فاصله آماری بین توزیع‌های خروجی را \textit{منقبض} کرده و آن‌ها را غیرقابل تشخیص می‌سازد.
			\pause
			\vspace{0.7cm}
			\item \textbf{پایگاه‌های داده مجاور:} دو پایگاه داده (مجموعه‌داده) $\mathcal{D}$ و $\mathcal{D}'$ مجاور هستند ($\mathcal{D} \sim \mathcal{D}'$) 
			اگر دقیقاً در یک رکورد تفاوت داشته باشند ($||\mathcal{D} - \mathcal{D}'||_1 \le 1$).
			
			\pause
			
			\vspace{0.5cm}
			مثال: یک مجموعه‌داده شامل قد افراد را در نظر بگیرید:
			
			\begin{itemize}
				\vspace{0.2cm}
				\item \textbf{پایگاه داده $\mathcal{D}_1$:} $\{170, 165, \mathbf{175}, 182\}$
				\item \textbf{پایگاه داده $\mathcal{D}_2$:} $\{170, 165, 182\}$
			\end{itemize}
		\end{itemize}
	\end{frame}
	
	\begin{frame}{محرمانگی تفاضلی: تعریف ریاضی}
		\begin{itemize}
			\item \textbf{تعریف رسمی ($(\varepsilon, \delta)$-DP):} یک مکانیزم تصادفی‌شده $\mathcal{M}: \mathcal{P}(\mathcal{X}) \rightarrow \mathcal{R}$ شرایط $(\varepsilon, \delta)$-محرمانگی تفاضلی را فراهم می‌کند، اگر برای تمام پایگاه‌های داده مجاور $\mathcal{D} \sim \mathcal{D}'$ \\ و تمام $S \subset \mathcal{R}$ داشته باشیم:
			\vspace{0.2cm}
			\begin{block}{}
				
				\begin{equation*}
					\mathbb{P}[\mathcal{M}(\mathcal{D}) \in S] \le e^\varepsilon \cdot \mathbb{P}[\mathcal{M}(\mathcal{D}') \in S] + \delta
				\end{equation*}
				
			\end{block}
			\pause
			\vspace{0.2cm}
			\item \textbf{بودجه محرمانگی ($\varepsilon \geq 0$):} پارامتری که میزان سخت‌گیری محرمانگی را کنترل می‌کند ($\varepsilon$ کوچکتر $\Rightarrow$ محرمانگی قوی‌تر، نویز بیشتر).
			\item \textbf{احتمال شکست ($\delta \in [0,1)$) :} احتمال بسیار ناچیزی که در آن تضمین سخت‌گیرانه $\varepsilon$-DP برقرار نباشد.
		\end{itemize}
	\end{frame}
	
	\begin{frame}{شهود هندسی: محرمانگی تفاضلی به عنوان انقباض اطلاعات}
		\begin{center}
			\begin{tikzpicture}[scale=0.9, transform shape]
				% Input Space (Datasets)
				\draw[thick, fill=blue!5, rounded corners=10pt] (-5,-2.5) rectangle (0,2.5);
				\node[font=\bfseries] at (-2.5, 2.8) {\rl{فضای ورودی $(\mathcal{X}^n)$}};
				
				\filldraw[black] (-2.5, 1.5) circle (2pt) node[above] {$\mathcal{D}$};
				\filldraw[black] (-2.5, -1.5) circle (2pt) node[below] {$\mathcal{D}'$};
				
				% Distance in input space
				\draw[<->, dashed, thick] (-2.5, 1.3) -- (-2.5, -1.3) node[midway, right] {$||\mathcal{D} - \mathcal{D}'||_1 = 1$};
				
				% Output Space (Probability Distributions)
				\draw[thick, fill=green!5, rounded corners=10pt] (2,-2.5) rectangle (7,2.5);
				\node[font=\bfseries] at (4.5, 2.8) {\rl{فضای خروجی (توزیع‌ها)}};
				
				% Distributions overlapping (Contraction)
				\draw[thick, blue, fill=blue!20, opacity=0.6] (4, 0) ellipse (1.5cm and 1cm);
				\node[blue!80!black, font=\bfseries] at (3.3, 0.3) {$\mathcal{M}(\mathcal{D})$};
				
				\draw[thick, red, fill=red!20, opacity=0.6] (5, 0) ellipse (1.5cm and 1cm);
				\node[red!80!black, font=\bfseries] at (5.7, -0.3) {$\mathcal{M}(\mathcal{D}')$};
				
				% Contracted Distance
				% \draw[<->, dashed, thick, draw=black!70] (4, -1.2) -- (5, -1.2) node[midway, below];
				% {$D_f \approx 0$};
				
				% Mapping arrows (Mechanism M)
				\draw[->, thick, draw=black!80, shorten >=5pt, shorten <=5pt] (-2.3, 1.5) to[out=20, in=150] node[above] {$\mathcal{M}$} (4, 1);
				\draw[->, thick, draw=black!80, shorten >=5pt, shorten <=5pt] (-2.3, -1.5) to[out=-20, in=-150] node[below] {$\mathcal{M}$} (5, -1);
				
				% Shrinking visual effect
				% \node[font=\bfseries\Large, text=orange!90!black] at (0, 0) {\textbf{انقباض}};
			\end{tikzpicture}
		\end{center}
		% \pause
		\vspace{0.3cm}
		\begin{itemize}
			% \item \textbf{بینش کلیدی:} مکانیزم تصادفی‌شده $\mathcal{M}$ به عنوان یک هسته مارکوف عمل می‌کند که فاصله آماری بین ورودی‌ها را \textit{منقبض} می‌کند (بسته به سطح محرمانگی)، و توزیع‌های خروجی آن‌ها را مجبور به هم‌پوشانی می‌کند.
			\item مکانیزم مانند فیلتری عمل می‌کند که فاصله آماری بین توزیع‌های خروجی را \textit{منقبض} کرده و آن‌ها را غیرقابل تشخیص می‌سازد.
		\end{itemize}
		
	\end{frame}
	
	\begin{frame}{مدل محلی: حذف متولی مورد اعتماد}
		\begin{itemize}
			\item \textbf{چالش محرمانگی تفاضلی مرکزی:}   
			\begin{itemize}
				\item برای جمع‌آوری داده‌های خام، کاملاً به یک \textit{متولی مورد اعتماد} وابسته است.
				\item یک نقطه شکست واحد ایجاد می‌کند (آسیب‌پذیر در برابر نشت داده‌ها، تهدیدات داخلی و احضاریه‌های قانونی).
			\end{itemize}
			
			\pause
			\vspace{0.4cm}
			\item \textbf{محرمانگی تفاضلی محلی \lr{(LDP)}:}
			\begin{itemize}
				\item مرز اعتماد را مستقیماً به دستگاه کاربر منتقل می‌کند.
				\item کاربران داده‌های خود را \textit{پیش از} ارسال به سرور، به‌صورت محلی آشفته می‌سازند.
			\end{itemize}
		\end{itemize}      
	\end{frame}
	
	
\end{document}